\begin{figure*}[p]
\centering

\begin{adjustbox}{max width=\linewidth, scale=0.85}
\begin{minipage}{\linewidth}
\begin{promptbox}[Turing Distinguishability Judge (Part 1/9)]
## Task

You are judging a pairwise Turing test for personalized user simulation.

Your task is to decide which candidate response was written by the real [HUMAN] user and which was written by an AI imitating that user.

## Inputs

You will receive:
1. Past responses from [HUMAN]
2. The current conversation context
3. Two candidate responses, Response A and Response B
4. An advisory watchlist for source-copy checks

One candidate is the real [HUMAN] response. The other candidate is AI-generated.

## User History

<|User History|>
{user_history}
<|End User History|>

## Context

<|Context|>
{context}
<|End Context|>

## Candidate Responses

### Response A

<|Response A|>
{response_a}
<|End Response A|>

### Response B

<|Response B|>
{response_b}
<|End Response B|>

## Advisory Watchlist

### Source-Copy Watchlist

<|Source-Copy Watchlist|>
{source_copy_watchlist}
<|End Source-Copy Watchlist|>


## Evaluation Procedure

Evaluate each response independently before comparing them. Reason backward from the response to the likely target, goal, and style.

Score each response on three criteria from 0.0 to 1.0.
\end{promptbox}
\end{minipage}
\end{adjustbox}

\vspace{-6pt}

\caption{Exact Turing distinguishability judge prompt used in the experiments, split across nine parts for typesetting.}
\label{fig:turing_prompt}
\end{figure*}

\begin{figure*}[p]
\centering

\begin{adjustbox}{max width=\linewidth, scale=0.85}
\begin{minipage}{\linewidth}
\begin{promptbox}[Turing Distinguishability Judge (Part 2/9)]
## Criteria

### 1. Immediate target

Identify the exact part of the current context that the response addresses.

Consider whether the response:
- Reacts to a specific point in the latest [OTHER] turn, the broader current context, or a plausible topic pivot
- Understands what the other person actually said
- Speaks from [HUMAN]'s perspective rather than [OTHER]'s perspective
- Is directed at [OTHER] or [OTHER - OP], rather than at someone merely described in the thread

A topic pivot can be valid when it is a plausible next move for [HUMAN]. Do not penalize a pivot only because it does not directly answer the previous turn, especially if [HUMAN] often pivots in the history or current context.

Do not reward broad persona fit if the response misstates a speaker's stance, role, or personal experience; assign a low immediate_target_score.

immediate_target_score:
- 0.0-0.2 = Wrong or absent target. The response does not reply to the current context, targets the wrong person, takes [OTHER]'s role, or responds to someone only described in the thread.
- 0.2-0.4 = Weak target. The response is broadly on topic but does not clearly address the latest [OTHER] turn, a relevant current-context point, or a plausible human topic pivot.
- 0.4-0.6 = Mixed target. The response addresses the general context or makes a possible pivot, but the exact target is ambiguous, unsupported, or only loosely connected to the current exchange.
- 0.6-0.8 = Good target. The response addresses a plausible part of the current context or makes a plausible topic pivot from [HUMAN]'s perspective, with only minor ambiguity.
- 0.8-1.0 = Strong target. The response clearly addresses the most plausible current-context point, or makes a strongly human-plausible pivot, from [HUMAN]'s perspective and directed at [OTHER] or [OTHER - OP].
\end{promptbox}
\end{minipage}
\end{adjustbox}
\end{figure*}

\begin{figure*}[p]
\centering

\begin{adjustbox}{max width=\linewidth, scale=0.85}
\begin{minipage}{\linewidth}
\begin{promptbox}[Turing Distinguishability Judge (Part 3/9)]
### 2. Human goal

Given the immediate target, identify what [HUMAN] was probably trying to do with the response. When judging the goals, use [HUMAN]'s history as evidence, not as a rigid script.

Consider whether the goal:
- Is plausible for [HUMAN] in this context, given [HUMAN]'s history
- Fits the local conversation pattern
- Preserves a specific personal or context-sensitive intent when one is available
- Avoids replacing [HUMAN]'s likely intent with a generic information request, assistant-like task, or narrator behavior

Do not reward a broadly useful or sensible goal if it replaces [HUMAN]'s likely local move, such as a joke, aside, anecdote, correction, quote-reply, agreement, disagreement, question, or brief reaction.

human_goal_score:
- 0.0-0.2 = Wrong or implausible goal. The response's goal does not fit [HUMAN], contradicts the context, invents a new task, or behaves like an assistant or narrator.
- 0.2-0.4 = Weak goal. The response has a generic or loosely plausible goal but mostly misses what [HUMAN] would likely try to do next, including any plausible pivot.
- 0.4-0.6 = Mixed goal. The response's goal is plausible by topic but changes, broadens, softens, escalates, or redirects [HUMAN]'s likely next move without strong support.
- 0.6-0.8 = Good goal. The response's goal, including a reasonable new-question or topic-pivot goal, is plausible for [HUMAN] in this context, with only minor uncertainty or drift.
- 0.8-1.0 = Strong goal. The response's goal naturally follows from [HUMAN]'s current intent, history, or local topic-switching pattern, while preserving personal or context-sensitive intent.
\end{promptbox}
\end{minipage}
\end{adjustbox}
\end{figure*}

\begin{figure*}[p]
\centering

\begin{adjustbox}{max width=\linewidth, scale=0.85}
\begin{minipage}{\linewidth}
\begin{promptbox}[Turing Distinguishability Judge (Part 4/9)]
### 3. Communication style and length

Judge how the response is written, separate from whether it is a question, statement, pivot, opinion, or personal reframing.

Compare the response to [HUMAN]'s past responses and the local framing. Consider:
- Wording and phrasing
- Tone, humor, bluntness, and emotion
- Length and level of detail
- Grammar, capitalization, punctuation, and misspellings
- Specificity and natural level of effort

Do not reward generic fluency. Smooth, polished, or conventionally well-written prose is not a style match unless [HUMAN] writes that way. Judge whether the response preserves [HUMAN]'s characteristic wording, rhythm, roughness, formatting, punctuation, humor, profanity, hedging, and level of elaboration.
Match length as a style feature. A response that is much shorter or much longer than [HUMAN]'s likely response should not receive a high communication_style_score unless [HUMAN]'s history supports that length in this local situation.
Clearly broken or artifact-like writing should dominate the judgment even when the target or goal seems plausible. Hard penalize artifact-like generations.

communication_style_score:
- 0.0-0.2 = Strong mismatch. Wording, rhythm, tone, length, grammar, punctuation, specificity, roughness, polish, or format strongly conflicts with [HUMAN]'s history and local framing, or the response is artifact-like.
- 0.2-0.4 = Weak style match. The response matches only easy surface cues, such as being short, informal, fluent, blunt, or emotional, but misses [HUMAN]'s distinctive wording, rhythm, roughness, punctuation, humor, specificity, or level of elaboration.
- 0.4-0.6 = Mixed style match. Some style cues fit, but important mismatches remain in length, polish, grammar, emotion, specificity, formatting, punctuation, or level of effort; smooth generic prose usually belongs here at best.
- 0.6-0.8 = Good style match. The response mostly matches [HUMAN]'s characteristic wording, rhythm, tone, length, formatting, punctuation, specificity, and level of elaboration, with only minor mismatch; length should be in the same rough range as [HUMAN]'s likely response.
- 0.8-1.0 = Strong style match. The response sounds specifically like [HUMAN], including distinctive phrasing, rhythm, roughness or polish, humor, profanity, hedging, formatting, punctuation, specificity, natural level of effort, and closely matched length.
\end{promptbox}
\end{minipage}
\end{adjustbox}
\end{figure*}

\begin{figure*}[p]
\centering

\begin{adjustbox}{max width=\linewidth, scale=0.85}
\begin{minipage}{\linewidth}
\begin{promptbox}[Turing Distinguishability Judge (Part 5/9)]
## Penalty Checks

Each penalty is scored from 0.0 to 1.0. Use the same scale for every penalty:
- 0.0 = No issue.
- 0.1-0.2 = Minor possible issue; mention it if relevant, but it should barely affect the judgment.
- 0.3-0.4 = Noticeable issue; the response is still plausibly human, but confidence should drop.
- 0.5-0.6 = Serious issue; the response has a substantial penalty-worthy flaw.
- 0.7-0.8 = Very strong issue; the response is unlikely to be human for this reason.
- 0.9-1.0 = Decisive issue; the response is almost certainly not human for this reason.

### Bad quote or source copy
The source-copy watchlist is an advisory 5-gram scan against user history and current context. Only assign source_copy_penalty for a response when the watchlist is triggered for that response. If the watchlist is off or not triggered for a response, assign 0.0 and do not invent source-copy violations from short phrases not shown in the watchlist.

When the watchlist is triggered, first decide whether the overlap is local uptake: the candidate may repeat words from the immediate context in order to quote, agree with, disagree with, answer, or otherwise react to that exact text. This is allowed even when the reused text is not explicitly marked with quotation marks, blockquote formatting, or attribution. Do not treat missing quotation marks, blockquote formatting, or attribution as source-copy evidence when the copied words come from the immediate context and the candidate is reacting to them. Do not penalize logical quotes without quotation marks. Do not penalize generic conversational frames, common question templates, pet phrases, or common sign-offs.

If copied text is not a natural quote or local uptake (i.e., a direct copy from the user history), assign a high source_copy_penalty. If the response quotes text but the quotation logic is confusing or contradictory, assign a high source_copy_penalty.
\end{promptbox}
\end{minipage}
\end{adjustbox}
\end{figure*}

\begin{figure*}[p]
\centering

\begin{adjustbox}{max width=\linewidth, scale=0.85}
\begin{minipage}{\linewidth}
\begin{promptbox}[Turing Distinguishability Judge (Part 6/9)]
### Wrong target or speaker role
Assign wrong_target_or_role_penalty when a response speaks from the wrong role, addresses the wrong speaker, or assigns an unsupported stance, experience, motive, relationship, or conflict to the speaker it addresses.

This includes responses that:
- take another speaker's role instead of [HUMAN]'s role, including by assigning [OTHER]'s or [OTHER - OP]'s views, experiences, motives, relationships, or conflicts to [HUMAN]
- address someone only described in the context instead of an actual current-context speaker, such as [OTHER] or [OTHER - OP]
- address a real current-context speaker but treat that speaker as holding a view they did not state
- assign the addressed speaker a personal experience, motive, relationship, or conflict absent from the exchange

Use a high wrong_target_or_role_penalty when the role or target error makes the response implausible as [HUMAN]'s reply.

### Unsupported adversarial reframing
Assign unsupported_adversarial_reframing_penalty when a response attacks, corrects, rebuts, or cynically reframes a claim that the current-context speakers did not actually make.

Broad persona fit alone is not enough. A response can sound like [HUMAN]'s general argumentative style while still targeting the wrong claim.

Watch for:
- generic reframes such as "the real issue is", "that's not X, that's Y", or "you're not X, you're Y"
- repeated questions, stacked sarcastic questions, or blunt challenges that do not logically follow from the exchange
- unsupported accusations combined with generic demands for evidence, validation, or clarification
- invented motives, conflicts, criticisms, or roles absent from the current context

When a response is framed as a pushback, rebuttal, or challenge, check whether the objection responds to the claim actually made. If it attacks a premise the other speaker did not rely on, misstates a speaker's stance, role, or personal experience, or reframes the issue into a different dispute, assign a high unsupported_adversarial_reframing_penalty.
\end{promptbox}
\end{minipage}
\end{adjustbox}
\end{figure*}

\begin{figure*}[p]
\centering

\begin{adjustbox}{max width=\linewidth, scale=0.85}
\begin{minipage}{\linewidth}
\begin{promptbox}[Turing Distinguishability Judge (Part 7/9)]
### Assistant-like response
Assign assistant_like_penalty when a response clearly reads like chatbot output rather than an organic human comment based on surface form. Watch for numbered or bulleted action plans, section headers, step-by-step framing, summary or conclusion sections, repeated direct-address coaching, generic reassurance, overgenerated transition phrases such as "But wait" or "But wait—", template phrases such as "You're absolutely right", "You're absolutely correct", "here's what you need to do", "let's break this down", "next steps", or "set boundaries", overuse of dashes "-", semicolons ";", or emojis, especially multiple emojis in a short response when not justified by [HUMAN]'s history or the local context, unnatural colloquial apostrophe abbreviations such as "'em" when not supported by [HUMAN]'s history or the local context, and other AI-written template patterns. Do not treat "[...quoted text...]" as assistant-like by itself; it may appear in human-written responses.

Do not penalize a response merely because it uses bullets, links, citations, dashes, or structured argumentation if that format is natural for the platform, thread type, or [HUMAN]'s history.

## Scoring and Rating

Compute:
base_score_a = immediate_target_score_a + human_goal_score_a + communication_style_score_a
base_score_b = immediate_target_score_b + human_goal_score_b + communication_style_score_b
penalty_a = ((source_copy_penalty_a + wrong_target_or_role_penalty_a + unsupported_adversarial_reframing_penalty_a + assistant_like_penalty_a) / 4) * 3
penalty_b = ((source_copy_penalty_b + wrong_target_or_role_penalty_b + unsupported_adversarial_reframing_penalty_b + assistant_like_penalty_b) / 4) * 3
response_a_score = max(0.0, base_score_a - penalty_a)
response_b_score = max(0.0, base_score_b - penalty_b)
score_gap = response_b_score - response_a_score

Convert score_gap to the final 1-7 rating:
- rating = 1 if score_gap <= -2.0
- rating = 2 if -2.0 < score_gap <= -1.0
- rating = 3 if -1.0 < score_gap <= -0.25
- rating = 4 if -0.25 < score_gap < 0.25
- rating = 5 if 0.25 <= score_gap < 1.0
- rating = 6 if 1.0 <= score_gap < 2.0
- rating = 7 if score_gap >= 2.0

Rating scale:
- 1 = Definitely A is the real human response
- 2 = Very likely A
- 3 = More likely A than B
- 4 = Cannot tell / equally likely
- 5 = More likely B than A
- 6 = Very likely B
- 7 = Definitely B is the real human response
\end{promptbox}
\end{minipage}
\end{adjustbox}
\end{figure*}

\begin{figure*}[p]
\centering

\begin{adjustbox}{max width=\linewidth, scale=0.85}
\begin{minipage}{\linewidth}
\begin{promptbox}[Turing Distinguishability Judge (Part 8/9)]
## Output Format

Return exactly one valid JSON object with this schema:
{
  "immediate_target_a": "<What exact part of the context is Response A reacting to, including any plausible topic pivot? Does it understand the latest turn, preserve the speakers' actual stance/role/personal experience, and avoid answering an invented position or conflict? Is it from [HUMAN]'s perspective and targeted to [OTHER] or [OTHER - OP]?>",
  "immediate_target_score_a": <number from 0.0 to 1.0>,
  "immediate_target_b": "<What exact part of the context is Response B reacting to, including any plausible topic pivot? Does it understand the latest turn, preserve the speakers' actual stance/role/personal experience, and avoid answering an invented position or conflict? Is it from [HUMAN]'s perspective and targeted to [OTHER] or [OTHER - OP]?>",
  "immediate_target_score_b": <number from 0.0 to 1.0>,
  "human_goal_a": "<Given the immediate target, what was [HUMAN] probably trying to do with Response A? Is that goal plausible for [HUMAN] in this context, specific or context-sensitive rather than only generic, and does it preserve [HUMAN]'s likely local move such as a joke, aside, anecdote, correction, quote-reply, agreement, disagreement, question, or brief reaction? Do not over-penalize a plausible pivot, blunt statement, opinion, or personal reframing.>",
  "human_goal_score_a": <number from 0.0 to 1.0>,
  "human_goal_b": "<Given the immediate target, what was [HUMAN] probably trying to do with Response B? Is that goal plausible for [HUMAN] in this context, specific or context-sensitive rather than only generic, and does it preserve [HUMAN]'s likely local move such as a joke, aside, anecdote, correction, quote-reply, agreement, disagreement, question, or brief reaction? Do not over-penalize a plausible pivot, blunt statement, opinion, or personal reframing.>",
  "human_goal_score_b": <number from 0.0 to 1.0>,
  "communication_style_a": "<Does Response A's wording, rhythm, tone, length, grammar, roughness or polish, humor, profanity, hedging, formatting, capitalization, punctuation, misspellings, specificity, and level of elaboration match [HUMAN]'s past responses and local framing? Do not reward generic fluency unless [HUMAN] writes that way; note whether length is in [HUMAN]'s likely range.>",
  "communication_style_score_a": <number from 0.0 to 1.0>,
  "communication_style_b": "<Does Response B's wording, rhythm, tone, length, grammar, roughness or polish, humor, profanity, hedging, formatting, capitalization, punctuation, misspellings, specificity, and level of elaboration match [HUMAN]'s past responses and local framing? Do not reward generic fluency unless [HUMAN] writes that way; note whether length is in [HUMAN]'s likely range.>",
  "communication_style_score_b": <number from 0.0 to 1.0>,
  "base_score_a": <immediate_target_score_a + human_goal_score_a + communication_style_score_a>,
  "base_score_b": <immediate_target_score_b + human_goal_score_b + communication_style_score_b>,
  "response_a_score": <number from 0.0 to 3.0>,
  "response_b_score": <number from 0.0 to 3.0>,
  "score_gap": <response_b_score - response_a_score>,
\end{promptbox}
\end{minipage}
\end{adjustbox}
\end{figure*}

\begin{figure*}[p]
\centering

\begin{adjustbox}{max width=\linewidth, scale=0.85}
\begin{minipage}{\linewidth}
\begin{promptbox}[Turing Distinguishability Judge (Part 9/9)]
  "response_a_source_copy": "<If source-copy watchlist is on for Response A, explain whether the matched text is a generic frame, common template, pet phrase, sign-off, natural quote, local uptake, direct copy from user history, or confusing/contradictory quote. If the watchlist is off, write an empty string.>",
  "source_copy_penalty_a": <number from 0.0 to 1.0>,
  "response_b_source_copy": "<If source-copy watchlist is on for Response B, explain whether the matched text is a generic frame, common template, pet phrase, sign-off, natural quote, local uptake, direct copy from user history, or confusing/contradictory quote. If the watchlist is off, write an empty string.>",
  "source_copy_penalty_b": <number from 0.0 to 1.0>,
  "response_a_wrong_target_or_role": "<Explain any wrong speaker role, wrong addressee, unsupported stance attribution, or unsupported personal experience/motive/relationship/conflict for Response A. If there is no issue, write an empty string.>",
  "wrong_target_or_role_penalty_a": <number from 0.0 to 1.0>,
  "response_b_wrong_target_or_role": "<Explain any wrong speaker role, wrong addressee, unsupported stance attribution, or unsupported personal experience/motive/relationship/conflict for Response B. If there is no issue, write an empty string.>",
  "wrong_target_or_role_penalty_b": <number from 0.0 to 1.0>,
  "response_a_unsupported_adversarial_reframing": "<Explain any unsupported attack, correction, cynical reframe, wrong-claim rebuttal, invented motive/conflict/criticism, or illogical challenge for Response A. If there is no issue, write an empty string.>",
  "unsupported_adversarial_reframing_penalty_a": <number from 0.0 to 1.0>,
  "response_b_unsupported_adversarial_reframing": "<Explain any unsupported attack, correction, cynical reframe, wrong-claim rebuttal, invented motive/conflict/criticism, or illogical challenge for Response B. If there is no issue, write an empty string.>",
  "unsupported_adversarial_reframing_penalty_b": <number from 0.0 to 1.0>,
  "response_a_assistant_like": "<Explain any assistant-like surface-form issue for Response A: generic chatbot formatting, template phrasing, repeated direct-address coaching, generic reassurance, overgenerated transition phrases, overuse of dashes, semicolons, emojis, unsupported colloquial apostrophe abbreviations, or artifact-like surface form. If there is no issue, write an empty string.>",
  "assistant_like_penalty_a": <number from 0.0 to 1.0>,
  "response_b_assistant_like": "<Explain any assistant-like surface-form issue for Response B: generic chatbot formatting, template phrasing, repeated direct-address coaching, generic reassurance, overgenerated transition phrases, overuse of dashes, semicolons, emojis, unsupported colloquial apostrophe abbreviations, or artifact-like surface form. If there is no issue, write an empty string.>",
  "assistant_like_penalty_b": <number from 0.0 to 1.0>,
  "penalty_a": <((source_copy_penalty_a + wrong_target_or_role_penalty_a + unsupported_adversarial_reframing_penalty_a + assistant_like_penalty_a) / 4) * 3>,
  "penalty_b": <((source_copy_penalty_b + wrong_target_or_role_penalty_b + unsupported_adversarial_reframing_penalty_b + assistant_like_penalty_b) / 4) * 3>,
  "reasoning": "<Concise explanation of the base scores, penalties, final score gap, topic-pivot evidence, goal specificity, style evidence, capitalization and punctuation evidence, and final rating.>",
  "rating": <integer from 1 to 7>
}

## Format Rules

- Output only valid JSON.
- Keep every explanation clear and concise.
- All criterion scores must be numbers from 0.0 to 1.0.
- All penalty scores must be numbers from 0.0 to 1.0.
- "base_score_a", "base_score_b", "penalty_a", and "penalty_b" must match the formulas above.
- "response_a_score" and "response_b_score" must be numbers from 0.0 to 3.0.
- "score_gap" must equal response_b_score - response_a_score.
- "rating" must be a single integer from 1 to 7.

Your output:
\end{promptbox}
\end{minipage}
\end{adjustbox}
\end{figure*}
